\documentclass[a4paper,openright,11pt]{book}
\usepackage[left=2.5cm,right=2.5cm,top=2.5cm,bottom=5cm]{geometry}
\usepackage[utf8]{inputenc}
\usepackage[spanish]{babel}
\usepackage{fancyhdr}
\pagestyle{fancy}
\usepackage{graphicx}
\usepackage{tikz}
\usepackage{venndiagram}
\usepackage{xcolor}
\usepackage{listings}
\usepackage{tabularx}
\usepackage{amsmath}
\usepackage{multicol}
\usepackage{amssymb}
\usepackage[utf8]{inputenc}
\headheight=98pt

% Cambie aquí la sangria a su preferencia 
\usepackage{parselines} 
\setlength{\parindent}{0em} 
\setlength{\parskip}{.0em}


% Cambie aquí la el estilo de númeración de secciones y subsecciones si lo desea
\renewcommand\thesection{\Roman{section}}
\renewcommand\thesubsection{\arabic{subsection}}

% Introduzca su información personal a continuación
\fancyhead[L]{\includegraphics[width=6cm]{logo-Konrad-1.png}}
\fancyhead[R]{Materia:\\
Código:\\
Nombre: 
\vspace{0.4cm}
}



\renewcommand{\headrulewidth}{1pt}
\renewcommand{\footrulewidth}{1pt}

\begin{document}
\include{Code}
% Documento hecho por Cameron Reyes
% Magnifico 
    \section{Puede crear una sección}
    \subsection{Puede crear una subsección}
        
        % Puede hacer matrices
        $$A=\begin{pmatrix} 
         a_{11}& a_{12} & \dotsc & a_{1n}\\ 
         a_{21}& a_{22} & \dotsc & a_{2n}\\
         \vdots & \vdots & \ddots&  \vdots\\
          a_{2m}& a_{2m} & \dotsc & a_{mn}
        \end{pmatrix} \hspace{1 cm} B=\begin{pmatrix} 
         b_{11}& b_{12} & \dotsc & b_{1n}\\ 
         b_{21}& b_{22} & \dotsc & b_{2n}\\
         \vdots & \vdots & \ddots&  \vdots\\
          b_{2m}& b_{2m} & \dotsc & b_{mn}
        \end{pmatrix} $$ \\
      
        %Creación de tablas
        \begin{center}
            \begin{tabular}{|c|c|}
               \hline Aquí una tabla & información \\ \hline
                      Información & información \\ \hline
            \end{tabular}
        \end{center}
        

        \subsection*{Puede hacer una subsección no enumerada}

        %Incluir ecuaciones 
        \begin{equation}
            \frac{|ax_1+b_1+c|}{\sqrt{a^2+b^2}}
        \end{equation}

        \subsection{Luego volver a enumerarlas}

        %Incluir graficos
        \begin{center}
            \includegraphics[scale =0.4]{logo-Konrad-1.png}
        \end{center}

        \section{Nuevas secciones}

        %incluir código
        \begin{lstlisting}
            print("Hola mundo")
        \end{lstlisting} \vspace{0.3 cm}
        
        %Puede Crear listas.
        \textbf{Lista:}
        \begin{itemize}
            \item Crear Listas
            \item Ejemplo
        \end{itemize}
    
\end{document}

